% TOP_PROBLEM: {"problemId": "problem.the-baum-connes-conjecture-for-the-k-theory-of-reduced-group-c-algebras", "canonicalTitle": "The Baum-Connes Conjecture for the K-Theory of Reduced Group C*-Algebras", "releaseRank": 53, "primaryDomain": "algebra_representation_category", "primaryDomainLabel": "Algebra, representation & category theory", "displayStatus": "open", "releaseStatus": "open"}
% !TeX program = lualatex
% Definitions and sources retained from research_batch_51_54.tex; research is in GPT_6_Astra_Ultra.
\documentclass[11pt]{article}
\usepackage[margin=25mm]{geometry}
\usepackage{fontspec}
\setmainfont{Latin Modern Roman}
\usepackage{amsmath,amssymb,mathtools}
\usepackage{unicode-math}
\setmathfont{Latin Modern Math}
\usepackage{enumitem,needspace}
\usepackage{xurl,hyperref}
\hypersetup{colorlinks=true,urlcolor=blue}
\setcounter{secnumdepth}{1}
\setlength{\parindent}{0pt}
\setlength{\parskip}{5pt}
\setlength{\emergencystretch}{3em}
\setlist[itemize]{leftmargin=3em,itemsep=5pt,topsep=4pt,parsep=0pt}
\allowdisplaybreaks
\newcommand{\N}{\mathbb{N}}
\newcommand{\Z}{\mathbb{Z}}
\newcommand{\Q}{\mathbb{Q}}
\newcommand{\R}{\mathbb{R}}
\newcommand{\C}{\mathbb{C}}
\newcommand{\F}{\mathbb{F}}
\newcommand{\A}{\mathbb{A}}
\newcommand{\PP}{\mathbb P}
\newcommand{\E}{\mathbb{E}}
\newcommand{\eps}{\varepsilon}
\newcommand{\dd}{\,\mathrm d}
\newcommand{\sref}[2]{\hyperlink{source:#1:#2}{[S#2]}}
\newcommand{\eref}[2]{\hyperlink{extra:#1:#2}{[E#2]}}
% Turn the next switch off for a shorter reading copy. The quotations preserve
% every exactTarget field, including qualifications omitted from a short summary.
\newif\ifshowcatalogue
\showcataloguetrue
\newcommand{\cataloguescope}[1]{\ifshowcatalogue
 \paragraph{Exact scope in the supplied catalogue.}
 \begin{quote}\small #1\end{quote}\fi}
\begin{document}
\setcounter{section}{52}
% ================================================================
% problemId: problem.the-baum-connes-conjecture-for-the-k-theory-of-reduced-group-c-algebras
% source releaseRank: 53; source releaseStatus: open
% exactTarget SHA-256: 1e80ca06bf1a687d4292bebe3976cb6de48e445e8334000523e6b6ad9407b631
\clearpage
\section{\texorpdfstring{The Baum-Connes Conjecture for the K-Theory of Reduced Group C*-Algebras}{The Baum-Connes Conjecture for the K-Theory of Reduced Group C*-Algebras}}\label{53}
\subsection{Definitions and mathematical statement}
Let $\underline EG$ be a classifying space for proper $G$-actions: fixed-point sets for compact subgroups are contractible and those for noncompact subgroups are empty. Topological $K$-theory here is the compact-support/direct-limit equivariant $K$-homology of $\underline EG$. The reduced group $C^*$-algebra $C_r^*(G)$ is the operator-norm completion of convolution operators in the left regular representation. The assembly map is
\[
 \mu_G:K_*^{\rm top}(G)\longrightarrow K_*(C_r^*(G)),\qquad *=0,1.
\]
The conjecture says it is an isomorphism for every second-countable locally compact $G$. This is the coefficient-free target; adding arbitrary coefficient algebras changes the conjecture.
\par\smallskip\textit{Formulation and concept references:} \sref{53}{1}.
\cataloguescope{Prove the Baum-Connes conjecture for every second-countable locally compact group G: the assembly map K\_top(G) \ensuremath{\to} K(C*\_r(G)) from the topological K-theory of the group, built from the classifying space for proper actions, to the K-theory of the reduced group C*-algebra is an isomorphism.}
\subsection{Short English statement}
Does geometric equivariant $K$-homology compute the operator-algebraic $K$-theory of every group in the stated class?
\subsection{Sources}
\begin{itemize}\raggedright
\item[\textnormal{[S1]}]\hypertarget{source:53:1}{} W. Lück, H. Reich, in Handbook of K-theory, Springer, 2005.
\url{https://arxiv.org/abs/math/0402405}.
{\small\textit{Catalogue formulation source.}}
\item[\textnormal{[S2]}]\hypertarget{source:53:2}{} There are no known counter-examples to the Baum-Connes conjecture.
\url{https://www.fields.utoronto.ca/talks/There-are-no-known-counter-examples-to-Baum-Connes-conjecture}.
{\small\textit{Catalogue research/status reference; not a proof endorsement.}}
\item[\textnormal{[S3]}]\hypertarget{source:53:3}{} Expanders and K-theory for group C* algebras.
\url{https://www.fields.utoronto.ca/talk-media/2/78/29/slides.pdf}.
{\small\textit{Catalogue research/status reference; not a proof endorsement.}}
\item[\textnormal{[S4]}]\hypertarget{source:53:4}{} Controlled Analytic Properties and the Quantitative Baum-Connes Conjecture.
\url{https://arxiv.org/abs/1908.02131}.
{\small\textit{Catalogue research/status reference; not a proof endorsement.}}
\item[\textnormal{[S5]}]\hypertarget{source:53:5}{} The Baum-Connes conjecture for extensions.
\url{https://arxiv.org/abs/2508.05726}.
{\small\textit{Catalogue research/status reference; not a proof endorsement.}}
\item[\textnormal{[E1]}]\hypertarget{extra:53:1}{} Shintaro Nishikawa, \emph{Crossed product approach to equivariant localization algebras}, arXiv:2108.11235v2, Theorem A, Definition 5.2 and Corollary 13.4; see also the introduction on gamma elements.
\url{https://arxiv.org/abs/2108.11235v2}.
{\small\textit{Added research source: Localization/crossed-product formulation for second-countable locally compact groups, including the kernel criterion and the semisplit evaluation sequence. Consulted 24 September 2026.}}
\item[\textnormal{[E2]}]\hypertarget{extra:53:2}{} Martin Finn-Sell, \emph{Controlled Analytic Properties and the Quantitative Baum--Connes Conjecture}, arXiv:1908.02131v2, author withdrawal notice dated 12 May 2026.
\url{https://arxiv.org/abs/1908.02131v2}.
{\small\textit{Added research source: Withdrawn: the notice reports a missing norm-control argument in the K-theory component. Cited for source status only, not as an established quantitative theorem. Consulted 24 September 2026.}}
\item[\textnormal{[E3]}]\hypertarget{extra:53:3}{} Ralf Meyer, \emph{The Baum--Connes conjecture for extensions}, arXiv:2508.05726v2, 29 August 2025, two-page note.
\url{https://arxiv.org/abs/2508.05726v2}.
{\small\textit{Added research source: Counterexample to an overstrong extension-permanence assertion with coefficients. This is not a counterexample to the coefficient-free target. Consulted 24 September 2026.}}
% END RESEARCH ADDITION REFERENCES-53
\end{itemize}

\end{document}
